% Generated draft: requires mathematical and source review.
\documentclass[11pt]{article}
\usepackage[margin=25mm]{geometry}
\usepackage{fontspec}
\setmainfont{Latin Modern Roman}
\usepackage{amsmath,amssymb,mathtools,unicode-math}
\setmathfont{Latin Modern Math}
\usepackage{enumitem,needspace,xurl,hyperref}
\setlength{\parindent}{0pt}
\setlength{\parskip}{5pt}
\setlength{\emergencystretch}{3em}
\newcommand{\N}{\mathbb{N}}
\newcommand{\Z}{\mathbb{Z}}
\newcommand{\Q}{\mathbb{Q}}
\newcommand{\R}{\mathbb{R}}
\newcommand{\C}{\mathbb{C}}
\newcommand{\F}{\mathbb{F}}
\newcommand{\A}{\mathbb{A}}
\newcommand{\PP}{\mathbb{P}}
\newcommand{\E}{\mathbb{E}}
\newcommand{\eps}{\varepsilon}
\newcommand{\dd}{\,\mathrm{d}}
\begin{document}
\section*{UnsolvedMath ID: 30003069}
\textbf{Source problem code:} catalogue entry
\subsection{Definitions and mathematical statement}
Let \(\operatorname{Gr}_{k,n}(\C)\) denote the Grassmannian of \(k\)-dimensional complex linear subspaces of \(\C^n\). The usual type \(A_{n-1}\) interpretation uses integers \(n\geq 2\), \(1\leq k<n\), and a positive integer \(r\), with \(\varpi_k\) the \(k\)-th fundamental weight of \(\mathfrak{sl}_n\). These parameter ranges are conventional interpretations; the supplied record does not explicitly specify them, particularly the domain of \(r\).

Under this interpretation, the Feigin--Fourier--Littelmann--Vinberg polytope \(\operatorname{FFLV}(r\varpi_k)\) can be defined using coordinates \(x_{a,b}\), \(1\leq a\leq b\leq n-1\). A Dyck path from \((i,i)\) to \((j,j)\), where \(i\leq j\), is a sequence of such pairs whose successive steps are \((a,b)\mapsto(a,b+1)\) or \((a,b)\mapsto(a+1,b)\). The polytope consists of the nonnegative coordinate vectors satisfying, for every such path \(D\),
\[
\sum_{(a,b)\in D}x_{a,b}\leq
\begin{cases}
r,&i\leq k\leq j,\\
0,&\text{otherwise}.
\end{cases}
\]
Coordinates outside \(a\leq k\leq b\) are forced to vanish; the remaining \(k(n-k)\) coordinates carry the standard integer lattice. This supplies a conventional type \(A\) definition, whose compatibility with the source's conventions requires review.

For a projective variety \(X\), a line bundle \(L\), a nonzero rational section \(\tau\) of \(L\), and a suitable full-rank valuation \(v\) of its rational function field with values in an ordered lattice \(\Z^d\), the Newton--Okounkov body is
\[
\Delta_v(X,L;\tau)=\overline{\operatorname{conv}}\left\{\frac{1}{m}v(s/\tau^m):m\in\Z_{>0},\ 0\ne s\in H^0(X,L^{\otimes m})\right\}\subseteq\R^d.
\]
The graph notation \(\operatorname{NO}_{G_0}^{r}\) refers to the source's associated Newton--Okounkov body. Interpreting it through this definition requires the exact graph, valuation, lattice, line bundle, and normalization. These choices are not specified in the supplied record; in particular, identifying its polarization with the \(r\)-th power of the Pl\"ucker line bundle remains a convention to check.

Two bodies are unimodularly equivalent if an affine isomorphism between their affine spans maps one body onto the other and identifies their induced integer lattices. In full-dimensional lattice coordinates this means a map \(x\mapsto Ax+b\), where \(A\in\operatorname{GL}_d(\Z)\) and \(b\in\Z^d\).

The source asks whether, for the distinguished \emph{non-mirror} plabic graph \(G_0\) associated with each intended \(\operatorname{Gr}_{k,n}(\C)\), and for every intended value of \(r\),
\[
\operatorname{NO}_{G_0}^{r}\ \text{is unimodularly equivalent to}\ \operatorname{FFLV}(r\varpi_k).
\]
The graph's adjective alone does not determine its construction. Consequently this is a conditional transcription of the intended question, pending recovery of the source's graph and conventions, rather than a fully specified standalone conjecture. The original statement omits the superscript \(r\) on the Newton--Okounkov body, whereas the supplied corrected statement includes it.

The supplied curation reports that the identification is known for \(\operatorname{Gr}_{3,6}\), and for the mirror or dual rectangle graph, while the question for the original \(G_0\) in general remains unresolved. These are supplied literature assessments, not independently verified mathematical conclusions; the mirror or dual graph must not be substituted for the stated non-mirror graph.
\subsection{Short English statement}
The question asks whether the Newton--Okounkov body attached to a particular non-mirror plabic graph for a complex Grassmannian is integrally affinely equivalent to the corresponding FFLV polytope. The graph, valuation conventions, and parameter range require source clarification. Supplied dated curation labels the problem partially solved and reports the case \(\operatorname{Gr}_{3,6}\) as known.
\subsection{Sources}
\begin{itemize}
\item Supplied local UnsolvedMath registry record, permanent ID \(30003069\), exact source problem code \texttt{catalogue entry}, title: Plabic Newton--Okounkov Bodies and FFLV Polytopes. Collection: Oberwolfach Reports Open Problems. The record gives proposed year 2016 and names no individual proposer.
\item Authoritative report cited by the record: \emph{Mini-Workshop: PBW Structures in Representation Theory} (2016), Oberwolfach Reports, pp.~615--651, DOI \texttt{10.4171/owr/2016/13}. \url{https://doi.org/10.4171/owr/2016/13}. The supplied background describes these reports as records of mathematical workshops and their open-problem discussions.
\item Supplied literature source, labeled ``plabic graph symmetries'': \url{https://doi.org/10.1016/j.crma.2018.05.003}.
\item Supplied literature source, labeled ``Newton--Okounkov bodies and Grassmannians'': \url{https://arxiv.org/abs/1712.00447}.
\item Literature assessment dated exactly \texttt{2026-08-21}: the record labels the problem \texttt{partially\_solved} and reports results for the mirror or dual rectangle graph and for \(\operatorname{Gr}_{3,6}\), while retaining the general question for the original \(G_0\) as unresolved. This is supplied dated curation derived from the OWR material, not independently verified here; its status and cited claims require independent verification.
\item The supplied statement status is \texttt{corrected\_verified}. Its verification note asserts comparison with the official report and preservation of the report's hypotheses, notation, and target. This assertion is part of the supplied record; no source visit or independent verification is claimed here.
\end{itemize}
\end{document}
